\chapter{扩散模型逆后验推导}
\label{app:diffusion-posterior}

在第 4.4 节中，我们指出真实逆后验
\[
q(\mathbf{x}_{t-1}\mid \mathbf{x}_t,\mathbf{x}_0)
\]
是高斯分布，并给出了其均值与方差的闭式表达式。这里给出的推导建立在与此前 PPCA 中相同的基本代数思想之上：一旦所有相关项都是高斯的，就可以把关于目标变量的二次项收集起来，再通过配方法得到结果。

\subsection*{步骤 1：从贝叶斯公式开始}

根据前向过程，有
\[
q(\mathbf{x}_{t-1}\mid \mathbf{x}_t,\mathbf{x}_0)
=
\frac{
q(\mathbf{x}_t\mid \mathbf{x}_{t-1},\mathbf{x}_0)\,
q(\mathbf{x}_{t-1}\mid \mathbf{x}_0)
}{
q(\mathbf{x}_t\mid \mathbf{x}_0)
}.
\]
由于分母不依赖于 \(\mathbf{x}_{t-1}\)，当把该式视为关于 \(\mathbf{x}_{t-1}\) 的函数时，分母只起归一化常数的作用。因此，在比例意义下，
\[
q(\mathbf{x}_{t-1}\mid \mathbf{x}_t,\mathbf{x}_0)
\propto
q(\mathbf{x}_t\mid \mathbf{x}_{t-1},\mathbf{x}_0)\,
q(\mathbf{x}_{t-1}\mid \mathbf{x}_0).
\]

由于前向过程满足 Markov 性，
\[
q(\mathbf{x}_t\mid \mathbf{x}_{t-1},\mathbf{x}_0)
=
q(\mathbf{x}_t\mid \mathbf{x}_{t-1}),
\]
因此两个高斯因子分别为
\[
q(\mathbf{x}_t\mid \mathbf{x}_{t-1})
=
\mathcal{N}\!\bigl(
\mathbf{x}_t\mid
\sqrt{\alpha_t}\,\mathbf{x}_{t-1},
(1-\alpha_t)\mathbf{I}
\bigr),
\]
以及
\[
q(\mathbf{x}_{t-1}\mid \mathbf{x}_0)
=
\mathcal{N}\!\bigl(
\mathbf{x}_{t-1}\mid
\sqrt{\bar{\alpha}_{t-1}}\,\mathbf{x}_0,
(1-\bar{\alpha}_{t-1})\mathbf{I}
\bigr).
\]

\subsection*{步骤 2：把乘积写成指数形式}

忽略与 \(\mathbf{x}_{t-1}\) 无关的归一化常数，有
\[
\begin{aligned}
q(\mathbf{x}_{t-1}\mid \mathbf{x}_t,\mathbf{x}_0)
&\propto
\exp\!\left(
-\frac{1}{2(1-\alpha_t)}
\left\|
\mathbf{x}_t-\sqrt{\alpha_t}\,\mathbf{x}_{t-1}
\right\|^2
\right)
\\[4pt]
&\qquad\times
\exp\!\left(
-\frac{1}{2(1-\bar{\alpha}_{t-1})}
\left\|
\mathbf{x}_{t-1}-\sqrt{\bar{\alpha}_{t-1}}\,\mathbf{x}_0
\right\|^2
\right).
\end{aligned}
\]
合并两个指数项得到
\[
\begin{aligned}
q(\mathbf{x}_{t-1}\mid \mathbf{x}_t,\mathbf{x}_0)
\propto
\exp\!\Bigg(
&-\frac{1}{2(1-\alpha_t)}
\left\|
\mathbf{x}_t-\sqrt{\alpha_t}\,\mathbf{x}_{t-1}
\right\|^2
\\
&-\frac{1}{2(1-\bar{\alpha}_{t-1})}
\left\|
\mathbf{x}_{t-1}-\sqrt{\bar{\alpha}_{t-1}}\,\mathbf{x}_0
\right\|^2
\Bigg).
\end{aligned}
\]

\subsection*{步骤 3：展开关于 \texorpdfstring{$\mathbf{x}_{t-1}$}{x\_\{t-1\}} 的二次项}

现在展开两个平方范数，只保留包含 \(\mathbf{x}_{t-1}\) 的项。

对于第一项，
\[
\begin{aligned}
\left\|
\mathbf{x}_t-\sqrt{\alpha_t}\,\mathbf{x}_{t-1}
\right\|^2
&=
\mathbf{x}_t^\top\mathbf{x}_t
-2\sqrt{\alpha_t}\,\mathbf{x}_t^\top \mathbf{x}_{t-1}
+\alpha_t\,\mathbf{x}_{t-1}^\top \mathbf{x}_{t-1}.
\end{aligned}
\]
对于第二项，
\[
\begin{aligned}
\left\|
\mathbf{x}_{t-1}-\sqrt{\bar{\alpha}_{t-1}}\,\mathbf{x}_0
\right\|^2
&=
\mathbf{x}_{t-1}^\top \mathbf{x}_{t-1}
-2\sqrt{\bar{\alpha}_{t-1}}\,\mathbf{x}_{t-1}^\top \mathbf{x}_0
+\bar{\alpha}_{t-1}\,\mathbf{x}_0^\top\mathbf{x}_0.
\end{aligned}
\]

将这些结果代入指数部分，并去掉与 \(\mathbf{x}_{t-1}\) 无关的项，得到
\[
\begin{aligned}
q(\mathbf{x}_{t-1}\mid \mathbf{x}_t,\mathbf{x}_0)
\propto
\exp\!\Bigg(
&-\frac{1}{2}
\left[
\frac{\alpha_t}{1-\alpha_t}
+
\frac{1}{1-\bar{\alpha}_{t-1}}
\right]
\mathbf{x}_{t-1}^\top \mathbf{x}_{t-1}
\\
&+
\left[
\frac{\sqrt{\alpha_t}}{1-\alpha_t}\,\mathbf{x}_t^\top
+
\frac{\sqrt{\bar{\alpha}_{t-1}}}{1-\bar{\alpha}_{t-1}}\,\mathbf{x}_0^\top
\right]
\mathbf{x}_{t-1}
\Bigg).
\end{aligned}
\]

此时它已经具有标准的高斯二次型形式
\[
-\frac12 A\,\mathbf{x}_{t-1}^\top \mathbf{x}_{t-1}
+
\mathbf{b}^\top \mathbf{x}_{t-1},
\]
其中
\[
A
=
\frac{\alpha_t}{1-\alpha_t}
+
\frac{1}{1-\bar{\alpha}_{t-1}},
\qquad
\mathbf{b}
=
\frac{\sqrt{\alpha_t}}{1-\alpha_t}\,\mathbf{x}_t
+
\frac{\sqrt{\bar{\alpha}_{t-1}}}{1-\bar{\alpha}_{t-1}}\,\mathbf{x}_0.
\]

\subsection*{步骤 4：配方}

利用标准恒等式
\[
-\frac12 A\,\mathbf{x}^\top\mathbf{x}
+
\mathbf{b}^\top\mathbf{x}
=
-\frac12 A
\left\|
\mathbf{x}-A^{-1}\mathbf{b}
\right\|^2
+\text{常数},
\]
可以看出，该后验是高斯分布，其协方差系数为
\[
\tilde{\beta}_t = A^{-1},
\]
均值为
\[
\tilde{\boldsymbol{\mu}}_t = A^{-1}\mathbf{b}.
\]

因此先化简 \(A\)：
\[
\begin{aligned}
A
&=
\frac{\alpha_t}{1-\alpha_t}
+
\frac{1}{1-\bar{\alpha}_{t-1}}
\\[4pt]
&=
\frac{
\alpha_t(1-\bar{\alpha}_{t-1}) + (1-\alpha_t)
}{
(1-\alpha_t)(1-\bar{\alpha}_{t-1})
}.
\end{aligned}
\]
由于 \(1-\alpha_t=\beta_t\)，且 \(\bar{\alpha}_t=\alpha_t\bar{\alpha}_{t-1}\)，分子变为
\[
\alpha_t - \alpha_t\bar{\alpha}_{t-1} + 1 - \alpha_t
=
1-\bar{\alpha}_t.
\]
因此
\[
A
=
\frac{
1-\bar{\alpha}_t
}{
\beta_t(1-\bar{\alpha}_{t-1})
},
\]
从而
\[
\boxed{
\tilde{\beta}_t
=
A^{-1}
=
\frac{1-\bar{\alpha}_{t-1}}{1-\bar{\alpha}_t}\,\beta_t.
}
\]

接下来计算均值：
\[
\tilde{\boldsymbol{\mu}}_t
=
A^{-1}\mathbf{b}
=
\tilde{\beta}_t
\left(
\frac{\sqrt{\alpha_t}}{1-\alpha_t}\,\mathbf{x}_t
+
\frac{\sqrt{\bar{\alpha}_{t-1}}}{1-\bar{\alpha}_{t-1}}\,\mathbf{x}_0
\right).
\]
代入
\(
\tilde{\beta}_t
=
\frac{1-\bar{\alpha}_{t-1}}{1-\bar{\alpha}_t}\beta_t
\)
以及 \(1-\alpha_t=\beta_t\)，可得
\[
\begin{aligned}
\tilde{\boldsymbol{\mu}}_t
&=
\frac{1-\bar{\alpha}_{t-1}}{1-\bar{\alpha}_t}\beta_t
\left(
\frac{\sqrt{\alpha_t}}{\beta_t}\,\mathbf{x}_t
+
\frac{\sqrt{\bar{\alpha}_{t-1}}}{1-\bar{\alpha}_{t-1}}\,\mathbf{x}_0
\right)
\\[6pt]
&=
\frac{\sqrt{\alpha_t}(1-\bar{\alpha}_{t-1})}{1-\bar{\alpha}_t}\,\mathbf{x}_t
+
\frac{\sqrt{\bar{\alpha}_{t-1}}\beta_t}{1-\bar{\alpha}_t}\,\mathbf{x}_0.
\end{aligned}
\]
因此
\[
\boxed{
\tilde{\boldsymbol{\mu}}_t(\mathbf{x}_t,\mathbf{x}_0)
=
\frac{
\sqrt{\bar{\alpha}_{t-1}}\,\beta_t
}{
1-\bar{\alpha}_t
}\,\mathbf{x}_0
+
\frac{
\sqrt{\alpha_t}(1-\bar{\alpha}_{t-1})
}{
1-\bar{\alpha}_t
}\,\mathbf{x}_t.
}
\]

\subsection*{步骤 5：最终结果}

综合以上各式，
\[
\boxed{
q(\mathbf{x}_{t-1}\mid \mathbf{x}_t,\mathbf{x}_0)
=
\mathcal{N}\!\Bigl(
\mathbf{x}_{t-1}\mid
\tilde{\boldsymbol{\mu}}_t(\mathbf{x}_t,\mathbf{x}_0),
\tilde{\beta}_t\mathbf{I}
\Bigr),
}
\]
其中
\[
\tilde{\boldsymbol{\mu}}_t(\mathbf{x}_t,\mathbf{x}_0)
=
\frac{
\sqrt{\bar{\alpha}_{t-1}}\,\beta_t
}{
1-\bar{\alpha}_t
}\,\mathbf{x}_0
+
\frac{
\sqrt{\alpha_t}(1-\bar{\alpha}_{t-1})
}{
1-\bar{\alpha}_t
}\,\mathbf{x}_t,
\]
并且
\[
\tilde{\beta}_t
=
\frac{
1-\bar{\alpha}_{t-1}
}{
1-\bar{\alpha}_t
}\,\beta_t.
\]

这就是 ELBO 中学习得到的去噪转移 \(p_\theta(\mathbf{x}_{t-1}\mid \mathbf{x}_t)\) 所需要拟合的精确局部逆向目标。
